% 原书分章PDF第14—18页：第30.6节全部12段，四因解释及全部22个哲学问句。
\section{因果性的一些根源与侧面}

% S30.6.P01
虽然真理是整个科学事业赖以建立的根本概念，但前面已经指出，它似乎不再在认识论或研究方法论中发挥关键作用。人们往往非正式地使用它，甚至不以助范畴词的方式（syncategorematically）引入它；它主要出现或现身于涉及较次要问题的争论。如前所述，这里将说明，真理如何在哲学与科学中另一个同样持久的主题语境里表现出来，即复杂的因果性概念，以及与之交织的因果作用、原因—结果关系、因果推断、因果推理或因果建模等概念。我们不打算依历史依据、约定式定义或方法论界线，试图区分这些概念，而是把因果性用作一个总括术语，涵盖上述概念群。物理必然性、逻辑必然性、不可避免性、充分性、决定论、规律性等概念，以及偶然性、概率或不规则性等反义概念，也构成因果性的内容，或在历史上与之相联。

% S30.6.P02
不过，与真理不同，因果性具有许多应用，并在科学哲学中呈现出新的研究方向；它在方法论研究中占据主导地位，在某种意义上甚至十分“时髦”，尤其是过去二十年来计算机科学与人工智能的发展使然。因果性似乎是一个古老而根本的概念，是“人是万物的尺度”（homo mensura）不可避免的一个方面。它首先似乎是一个直觉概念、常识概念或自然范畴（按亚里士多德或康德意义上的用法）；我们不断使用它，以便在生存斗争中适应环境，解释自己的经验，理解自身、自己在宇宙中的位置，以及人类存在的偶然境遇。因此，因果性在日常语言使用中无处不在。它明确地包含在许多连接词和众多“因果”动词之中，包括造成（causing）、诱发（inducing）、影响（affecting）、触发（triggering）、启动（initiating）、导致（leading to）、使发生（effecting）、施加影响（influencing）、产生（producing）、结果成为（resulting in），等等。在某种意义上，所有及物动词本质上都具有因果性。其中一些表达物理必然性或偶然性，另一些则预设一个主体或行动者，由其有意而为地使一个客体发生变化。这样，早在哲学或科学曙光出现以前，因果性便已经显现出来。

% S30.6.P03
从那时起，这个概念便呈现出多种面貌、内涵和解释，并在数百年间不断演变。它已表明自己具有变色龙般的特性：适应其被运用或发展的哲学语境或理论。正如将要看到的，它似乎高度依赖语境，因此与逻辑必然性、物理必然性、偶然性、巧合、机遇、目的论和决定论联系在一起；在实在论、唯心论、经验主义、工具主义或理性主义等不同立场中，也呈现出不同面貌。尽管概念变化从未停止，各个时代对“因果性”这个术语的使用却具有显著的一致性。

% S30.6.P04
更重要的是，直到今天，它仍在科学哲学、统计学和数据科学中引发混乱与争论。首先，即便最基本、最简洁的表达“$X$导致$Y$”（X causes Y），似乎也会打开一个装满问题与微妙之处的潘多拉魔盒。这个概念涵盖十分不同的思想与方法。我们可以看到决定论与概率论立场、差异制造进路、基于规律性的方法、反事实进路、干预主义进路、机制理论、类型与个例的两难问题、原因与理由之间被认为存在的鸿沟、个体或群体层面的说明、物理或心灵主义解释、单因与多因、定性或定量说明、离散与连续过程中的表现、观察与实验情境中的表现，以及直接与间接因果性，等等。这些进路或解释显然既非相互排斥，也非完全穷尽；这个清单很容易继续扩展（Tacq，2011；Williamson，2009）。一旦脱离其历史和哲学背景，它们便呈现出万花筒般的图景，只容许一种多元而不统一的说明，尤其考虑到前面勾勒的那幅令人遗憾的三联画时更是如此。几个历史细节或许能略微澄清这个问题。

% S30.6.P05
英国哲学家、数学家阿尔弗雷德·诺思·怀特海（Alfred North Whitehead，1861—1947）至今仍受到哲学家的推崇，因为他批评科学中的唯物主义，尤其因为他是过程哲学的倡导者与先驱。西方传统深受巴门尼德著作、柏拉图超越时空的理念或形式，以及亚里士多德的实体思想与本质主义的影响；其中，巴门尼德把存在的不变性确立为一项本体论原则。怀特海表达了不同的看法。他更接近于追随以弗所的赫拉克利特这位在“过程哲学”一词出现之前的过程哲学家的足迹，发展出一种“有机体哲学”，把实在理解为由连续的“现实机缘”（actual occasions）构成的宏大过程。“物质”对象是这些当下事件之中的实际模式；这些事件既可以是客观的，也可以是主观的，从其他事件获取信息，并向其他事件传递信息。此外，经验与创造性都是起主导作用的构成性原则。由此，这些事件也具有因果活动性，使整个世界事件成为一个连续过程；变异与变化是首要的，并非从静态状态派生而来，也不能还原为静态状态。理解实在，意味着理解过程、程序和机制；这些概念在当代因果性思想中仍可辨认出来。

% S30.6.P06
由于结合了这些思想，而且不再以贬义方式看待变异与变化，怀特海与许多不同的思想家和同时代人联系起来，包括法国哲学家亨利·柏格森、生物学家华莱士与达尔文，以及统计学家高尔顿与皮尔逊（Starmans，2011a）。另一方面，今天的逻辑学领域主要把怀特海视为伯特兰·罗素宏著《数学原理》（Principia Mathematica）的合著者。该书于1910—1913年间出版，构成数学哲学中逻辑立场的巅峰。

% S30.6.P07
虽然怀特海的多才多艺与思想深度在这里并非题内话（out of order），但具有讽刺意味的是，他被引用最多的那句话，却完全属于另一种性质。他在1929年的《过程与实在》（Process and Reality）中断言：“对欧洲哲学传统最稳妥的一般概括，就是它由一连串对柏拉图的脚注组成。”人们可以把这个大胆判断理解为论战性的讥讽，也可以把它视为一个笑话，强调柏拉图哲学的永恒性。人们也可以试着从西方哲学中找出一组重大且至今仍有现实意义的经典主题，然后检视它们在多大程度上植根于柏拉图的著作，或受到其启发。进行这样的考察，就不可避免地会遇到持久而复杂的因果性问题。尽管原子论者和斯多亚学派已经探索过必然性与决定论概念，仍可以说，正是柏拉图在其对话《斐多篇》（Phaedo）中，通过若干极具问题性的片段，纳入了因果性理论的开端。

% S30.6.P08
另一方面，他的学生亚里士多德在《后分析篇》（Analytica Posteriora）和《物理学》（Physics）中，以四因（causae，亦即 aitiai）学说，提供了系统得多、也更有影响力的论述。认识一个事物，意味着认识其 aitia。要完整解释和理解某一特定实体，就必须回答涉及四因的四个问题：质料因（causa materialis）——“它由什么制成”；动力因（causa efficiens）——“谁制造了它，或发起了制造”；目的因（causa finalis）——“它是为了什么而被制造”；形式因（causa formalis）——“是什么使这个实体成为它自身，构成它的本质，而非其他东西”。这套学说处于其形而上学和认识论的核心，直到十七世纪，都在神学、哲学和科学中引发许多解释与修正，使亚里士多德成为因果性理论的奠基者。

% S30.6.P09
此后，围绕因果性的问题一直主导形而上学和认识论，直至十九世纪；这些问题可见于阿奎那（Aquinas）、培根、笛卡尔、伽利略（Galilei）、斯宾诺莎、霍布斯、牛顿、莱布尼茨、休谟、康德、斯图亚特·密尔和 C.S. Peirce 等不同思想家的著作。上述所有思想家以及更多人的著作，都表现出这个概念对语境的依赖和变色龙般的特征：有时是明确的，即他们对这个概念进行理论思考；有时是隐含的，即他们把它作为原始术语来使用，或以其他方式视为理所当然。事实上，因果性一次又一次地被重新发明，受到相关思想家的形而上学与认识论观点的影响，并与之紧密相联。由于哲学中的理论或观点很少被证伪和取代，许多相互冲突的思想延续下来，演变为当今认识论中关于因果性的多元观点。确实，即使在今天，哲学家仍在仔细考察因果关系的本体论与认识论方面，以及要作出有意义的因果陈述和因果推理，需要哪些假设与其他条件。如前所述，因果性、科学陈述与规律／自然规律之间的互动和相互依赖，仍构成一般科学哲学的一个重要领域或焦点。

% S30.6.P10
承认思想史中形成因果性概念的非凡演变，但绝不自称要作详细的历史综述，我们开始概念分析时，多少还是要把这个概念拉入经典的“哲学三角”：实在、心灵／观念和语言，以及这些概念之间微妙的相互作用。这里提出的问题，在过去两千年里一直主导着西方形而上学与认识论。是否存在一个客观、独立的世界（实在论），还是说，实在由心灵构建而成（唯心论、建构主义）？这个世界是否可知，如果可知，又如何认识？我们能否对这个世界作出真陈述，又如何表征它们？真理又是什么？

% S30.6.P11
聚焦于因果性与原因—结果关系，可以把这些问题具体化如下。从实在出发：因果性是否存在；如果存在，它是以静态方式（作为一种状态、事实或事实情境）存在，还是以动态方式（作为具有生产性质的物理过程）存在？它是否是由某个因果行动者发起的行动，该行动者从某种立场出发有意地实施这一行动？应不应该把这个过程看作一连串不间断的步骤，一种通向最终结果的原因链？这个过程是必然的、决定性的，还是随机的，亦即不确定的？它是连续发生，还是离散地／分步发生？这个过程是否还需要一种作用机制？能否存在多个原因，它们能否相互影响（加强、颠覆或抵消）？为了确立因果性，我们应该检视实在的一部分，还是实施干预？

% S30.6.P12
从心灵出发：因果性是否是心灵的一种建构，是把经验与期待投射到实在之上（休谟），还是人类能力的一项特征以及感知的条件（康德），从而预设真实的本体世界与现象世界的二元论？我们能否凭借感官与认知来感知它、假设它或认识它？再从语言出发：我们能否用自然语言、逻辑形式体系（时态逻辑、过程逻辑、行动逻辑、情境演算），或数学语言来表征它或将其操作化；若能使用数学语言，又应使用确定性函数、（线性）方程组，还是概率，等等？概率是否足够，还是我们需要 Pearl 的 do 演算这样的语言？我们是否需要推断统计学来表征和理解它；如果需要，又应如何解释结果？是借助可能世界语义学，还是通过反事实？我们是否应该量化、测量，然后确定平均值，把它作为因果性或因果效应的操作化方式（正如潜在结局进路所假定的那样）？我们能否用必要条件和充分条件（如 INUS 条件）来指出它；如果可以，这些条件能否有限公理化，或以一个有限清单被把握？在这样做时，我们能否避开人工智能领域中声名狼藉的框架问题、量化问题（quantification problem）或衍生问题（ramification problem）？能否存在一个统一的因果性理论，整合许多不同的方面与内涵？多种立场有若干历史根源，而且似乎不利于这个概念在研究方法论中的进一步发展，下一节将对此作出概述。
